% 本文件由 LaTeX 排版 Agent 根据有效 translation.json 自动生成。
% author=Augustine of Hippo; work_id=1306; title=论相信的益处

\chapter{论相信的益处}

% structure: p0001 source=h1 target=omitted reason=page-title-already-rendered-by-parent

\Emph{\WorkTitle{Retract.}} i. cap. 14. 此外，如今在希波勒及，我作为司铎写了一本书\WorkTitle{论相信的益处}，给一位被摩尼教徒欺骗的朋友，我知道他仍执迷于该谬误，并嘲笑天主教信仰学派，因人们被要求相信，却没有通过最确定的方法被教导什么是真理。在这本书中，我说了等等。***。这本书的开头是：\Emph{\Quote{Si mihi Honorate, unum atque idem videretur esse.}}\par

圣奥斯定将\WorkTitle{论相信的益处}列为他任司铎时所写的第一本书，他于公元391年初在希波被擢升为司铎。他为之写作的那个人曾被自己引入歧途，似乎后来被挽回，至少如果他就是公元412年左右从迦太基写信给圣奥斯定、提出数个问题且圣奥斯定回信（第140封书信）的那一位。卡西奥多罗斯称他为司铎，尽管当时他尚未受洗。在第83封书信中，圣奥斯定提到另一位名叫奥诺拉特的司铎之死。晚年，他还写了第228封书信给一位同名的塔本纳主教。——（本笃会版）\par

\WorkTitle{订正录}中的评论已附于相应段落的注释中。\par

1. 奥诺拉特，若异端者与相信异端者在我看来是同一回事，我就应判断自己有义务在此事上保持口舌和笔墨的沉默。然而，如今这两者之间有天壤之别：因为我认为，异端者是出于某种暂时的利益，尤其是为了自己的荣耀和优越，而创立或追随虚假的新奇见解的人；但相信此类人者，是因某种真理与虔敬的想象而被欺骗的人。既然如此，我认为自己有义务不向你隐瞒我对寻求和持守真理的看法——你知道，我们从最年轻时就以极大的爱火热爱这真理；但这远离虚妄之人的心智，他们过于深入并陷于这些物质事物，以为除了通过那五个众所周知的肉体报告者所感知的东西外别无他物；他们将由此获得的印象和形象随身携带，即使试图脱离感官时也是如此；并以这致命而最欺骗性的规则，自以为最正确地衡量了真理那不可言说的深处。我至亲爱的朋友，一个人不仅说而且认为自己已发现了真理，这是最容易不过的事；但事实上这有多么困难，我希望你从此信中可以察觉。为使这封信对你有益，或至少无害，也对所有偶然读到它的人有益，我已向天主祈祷，并仍在祈祷；我希望它如此，因为我深知自己是以虔敬和友爱之心、而非为虚名或轻浮的炫耀而执笔。\par

2. 我的目的，若可能，是向你证明：摩尼教徒亵渎而鲁莽地抨击那些人，即那些跟随天主教信仰的权威、在能够凝视那纯洁心灵所领悟的真理之前，藉相信而预先武装自己、为即将光照他们的天主预备自己的人。因为你知道，奥诺拉特，我们落入这样的人手中，别无其他原因，就是因为他们曾说：他们能抛开一切权威的恐惧，以纯粹而简单的理性，将那些愿意聆听他们的人领到天主面前，并从一切错误中解放出来。因为，是什么别的理由，在我几乎九年的时间里，摒弃我自幼由父母所建立的宗教，成为那些人的追随者和勤勉听众呢？不就是因为他们说我们被迷信所恐吓，被命令在理性之前相信，而他们则不强迫任何人相信，除非先讨论并阐明真理？谁能不被这样的承诺所诱惑呢？尤其是渴望真理的年轻心灵，而且是一颗在学堂里因某些博学之人的讨论而变得骄傲而多言的心。他们当时遇见我这样的人，我确实蔑视如老妇之寓言，渴望把握并畅饮他们所承诺的开放而纯粹的真理。但是，另一方面，是什么理由阻止了我完全加入他们，使我停留在他们所谓的\Quote{听者}等级中，以致没有放弃此世的希望与事业呢？除非我注意到，他们更善于雄辩而充分地驳斥他人，却在自己的主张证明上不够坚定可靠。但关于我自己，我又能说什么呢？我已是天主教基督徒。如今，在长久的干渴之后，我几乎耗尽且干涸的乳头，我以全部贪婪回归了它们，并以更深的哭泣和叹息将它们更深地摇动并拧挤出来，以便流出足够的东西来滋润我这样受影响的人，并带回生命与安全的希望。那么，关于我自己，我该说什么呢？你，尚未成为基督徒，通过我的鼓励，虽然你大大憎恶他们，仍勉强被说服应当听他们并考验他们，请问，是什么别的使你喜悦呢？我恳求你回想，难道不是一种对理由的巨大自信和许诺吗？因为他们以极大的丰富和激烈长篇大论地争论无知之人的错误——这件事我后来才学到，对任何受过中等教育的人来说都最容易不过；即使他们将自己的一些东西植入我们心中，我们认为必须保留它，因为我们没有遇到其他可以同意的东西。他们对我们所做的，正是狡猾的捕鸟者所惯做的：他们在水边涂有粘鸟胶的枝条来欺骗口渴的鸟。因为他们用某种方式填满并遮盖周围的其他水域，或用令人惊慌的手段把它们吓走，使它们不是出于选择，而是出于缺乏而落入他们的网罗。\par

3. 但我为何不回答自己说，这些优美而巧妙的比喻以及此类指责，任何敌人都可以极尽才智与讽刺地向一切事物的教导者倾倒出来？然而我以为，应当在信中插入一些这类内容，以便劝告他们停止这种行径；这样，正如他所说，撇开老生常谈的琐事，可以让实质与实质较量，理由与理由辩驳。因此，让他们放弃那句常挂在嘴边、仿佛非说不可的话吧，当某个长期听讲的人离开他们时，他们就说：\Quote{光已在他身上穿过。}因为您，我主要关心的人，（因为我并不太过担忧他们，）看到了这话是何等空洞，任何人都极易挑剔。所以此事我留待您的智慧去考虑。因为我不担心您会认为我充满了内在的光，当我被这世上生活纠缠时，怀有昏暗的希望：娇妻之美、财富之奢、荣华之虚，以及其他一切有害致命的享乐。这一切，您并非不知道，我在作他们认真听讲者的时候，从未停止过渴望与希望。我并非归咎于他们的教导；因为我也承认，他们同样谨慎地劝人远离这些。但现在说当我从这一切事物的阴影中转身，决心仅以满足身体健康的饮食为满足时，却被光抛弃；而在我喜爱并沉浸于这些事物时，却是光明照耀、闪闪发亮——这样说的人，用最温和的话说，对于那些他喜欢长篇大论的事情，缺乏敏锐的洞察力。但，如果您愿意，让我们回到手头的问题吧。\par

4. 因为您很清楚，摩尼教徒通过指责天主教信仰，尤其是通过撕裂并毁谤\WorkTitle{旧约}来打动无知者；而他们完全不知道，这些事应当如何理解，它们以何种延伸方式有益地进入那些尚只能哭泣的灵魂的血脉与骨髓之中。因为在它们中有一些事情，对于那些愚昧且不顾惜自己心灵的人（这样的人很多）会造成一些轻微冒犯，所以它们可以受到大众化的指责；但由于其中所包含的奥秘，它们不能由许多人在大众化方式下得到辩护。然而那些知道如何辩护的少数人并不喜欢公开且多受议论的争辩，因此他们对那些最热切寻求他们的人之外的人所知甚少。因此，关于摩尼教徒这种鲁莽行为——他们借此指责\WorkTitle{旧约}和天主教信仰——我恳求您倾听那些感动我的思考。但我渴望并希望您能以我说这些话时的同样精神来接受它们。因为天主，祂知道我的良心秘密，知道我在这篇讲话中绝无任何邪恶诡诈；相反，正如我认为应当被接受的那样，是为了证明真理，而这正是我们久已决意为之生活的一件事；并且怀着难以置信的忧虑，惟恐我与您一同犯错可能非常容易，而与您一同坚持正道则最为困难（不用更强硬的词）。但我大胆预期，在这希望中——我希望您将与我们一同持守智慧之路——祂不会舍弃我，这位我已奉献给祂、日夜努力仰望的天主；并且由于我的罪过和过去的习性，我的心眼被软弱意见的打击所伤，我知道自己软弱无力，常常流泪哀求，并且如同在长期盲目与黑暗之后，眼睛勉强睁开，却仍因频繁的跳动与回避而拒绝它们所渴望的光；特别当有人努力向它们展示太阳本身时；现在正是我的情形——我不否认灵魂有一种不可言说且独特的美善，是心灵所看见的；并且我含着眼泪和呻吟承认自己尚不配得它。那么，如果我毫不伪装，如果我被责任所引导，如果我热爱真理，如果我珍视友谊，如果我非常担心您受欺骗，祂就不会舍弃我。\par

5. 因此，所有被称为\WorkTitle{旧约}的圣经，以四重方式传授给那些渴望认识它的人：依照\OriginalTerm{历史}{history}、依照\OriginalTerm{原因论}{ætiology}、依照\OriginalTerm{类比}{analogy}、依照\OriginalTerm{寓意}{allegory}。不要因为我使用希腊词语就认为我愚蠢。首先，因为我如此领受了，我不敢以别于我所领受的方式向您表明。其次，您自己明白，我们并没有使用这类术语的习惯；如果我加以翻译并造出新词，那才更愚蠢；但如果我用迂回的说法，就会在论述中失去自由。我只恳求您相信，无论我在何处犯错，都不会在各事上骄傲自大。因此，举例来说，依照历史传授，即教导所写之事或所行之事，或未行但仅以行文记下之事。依照原因论，即说明某件事或某句话为何发生或说出。依照类比，即表明\WorkTitle{旧约}与\WorkTitle{新约}两约互不矛盾。依照寓意，即教导某些所写之事不应仅按字面理解，而应以其象征意义来领会。\par

6. 所有这些方式，我们的主耶稣基督和祂的宗徒们都使用过。因为当有人反对祂的门徒在安息日掐麦穗时，祂引用了历史实例：祂说：\Quote{你们没有读过\Person{达味}在饥饿时和随从者所做的事吗？他怎样进了天主的殿，吃了供饼，这饼只有司铎才能吃，他和随从者都不可以吃？}但这个实例属于\OriginalTerm{原由}{\GreekText{αἰτιολογία}}，即当基督禁止休妻（除非为淫乱的缘故），而询问祂的人引证\Person{梅瑟}在给予休书后许可了休妻，祂说：\Quote{\Person{梅瑟}是因为你们的心硬才准许了这事。}这里给出了理由，说明为什么\Person{梅瑟}一时准许了这事；而基督的这条诫命似乎表明如今时代不同了。但要解释这些时代的变化，以及它们由天主眷顾的某种奇妙安排所制定和确立的秩序，说来话长。\par

7. 再者，类比——藉此可以清楚看到两约的一致性——我为什么要说所有人都使用了它，而他们又屈服于谁的权威呢？然而他们自己能够思考：他们惯常所说的有多少东西是某些（我不知道是谁的）真理败坏者插入圣经的呢？我经常认为他们的这种说法非常软弱，即使在我还是他们的听众时也是如此；不单是我，还有你们（我记得很清楚），以及我们所有人，都试图在形成判断时比一般听众更加谨慎。但现在，许多过去常常动摇我的事情已经被解释并清楚显明给我了；我指的是那些事情，他们的言论在大部分情况下自我夸耀，并且在没有对手的情况下越是自由、越是安全地高谈阔论；似乎在我看来，没有比他们断言圣经被腐化更无耻的了，或者（用更温和的说法）更粗心和无力的了；然而，在如此近的事情上，并没有现存的抄本可以证明它。因为如果他们断言，他们认为不应该完全接受它们，因为它们是由那些他们认为没有写出真理的人所写的；无论如何，他们的拒绝会更正确，或者他们的错误会更自然。因为这就是他们在名为\WorkTitle{宗徒大事录}的书上所做的事。当我独自思考他们的这个诡计时，我不禁惊讶。因为我抱怨的不是这些人的智慧不足，而是缺乏普通理解力。因为那本书有如此重大的事情，类似于他们所接受的，以至于我认为拒绝接受这本书也是极大的愚蠢，并且如果有什么事情冒犯了他们，就称之为虚假和插入。或者，如果这种语言是无耻的（事实上就是如此），那么为什么在\Person{保禄}的书信、在四福音书中，他们认为它们有何价值呢？在这些书中，我确信按比例来说，有更多的事情比在那本书中可能有的更多，而他们却认为这些是被伪造者插入的。但这正是我所认为的情况，我请求你们与我一起以尽可能平静和清晰的判断来思考它。因为你们知道，他们试图将他们创始人\Person{摩尼}的人格纳入宗徒的行列，他们说主应许要派遣给祂门徒的圣神，已经通过他来到我们这里。因此，如果他们接受了那卷明确记载圣神降临的\WorkTitle{宗徒大事录}，他们就无法说它是被篡改的。因为他们坚持认为在\Person{摩尼}本人时代之前，就有某些（我不知道是谁的）神圣书籍的伪造者；而且这些书被那些希望将犹太人的法律与福音结合起来的人伪造了。但关于圣神，他们不能说这样的话，除非他们断言那些人预见到了将来会出现的\Person{摩尼}，并说圣神是通过他派遣的，并在他们的书中预先写下了反对\Person{摩尼}的话。但关于圣神，我们将在别处更清楚地讨论。现在让我们回到我的目的。\par

8. 因为旧约中的历史、\OriginalTerm{原因论}{ætiology}、\OriginalTerm{类比}{analogy}和\OriginalTerm{寓意}{allegory}在新约中都已找到，我认为这已经充分证明了；现在需要展示寓意的部分。我们的救赎主亲口在福音中使用了旧约的寓意。祂说：\BibleQuote{这一代寻求征兆，除了约纳先知的征兆外，不会给它们任何征兆。正如约纳三天三夜在大鱼腹中，人子也要三天三夜在地心里。}因为为什么我要提及宗徒\Person{保禄}，他在致\WorkTitle{格林多前书}中表明，就连出谷的历史也是未来基督徒子民的寓意。\BibleQuote{但我不愿你们不知道，弟兄们，我们的祖先都曾在云下，都从海中经过，都曾在云中和海中受洗归于\Person{梅瑟}，都吃过同样的神粮，都喝过同样的神饮；因为他们所饮的，是来自伴随他们的神磐石，那磐石就是\Person{基督}。但在他们中大多数人，天主并不喜悦，因为他们倒毙在旷野里。这些事发生在我们身上，作为预像，叫我们不要贪恋恶事，如同他们贪恋过一样。我们也不可拜偶像，像他们中有些人一样，正如经上所记载：\InnerQuote{百姓坐下吃喝，起来玩耍。}我们也不可淫乱，像他们中有些人淫乱，一天就倒毙了两万三千人。我们也不可试探基督，像他们中有些人试探过，就被蛇所灭。我们也不可抱怨，像他们中有些人抱怨过，就被毁灭者所灭。这一切事发生，都是他们的预像，并且记录下来，是为了警戒我们这些末世的人。}宗徒中也有一种寓意，它确实与当前的问题密切相关，原因是他们自己经常提出它并在辩论中展示它。因为同一\Person{保禄}在致\WorkTitle{迦拉达书}中说：\BibleQuote{因为经上记载：\InnerQuote{\Person{亚巴郎}有两个儿子，一个由婢女所生，一个由自由妇人所生。那由婢女所生的，是按血肉生的；那由自由妇人所生的，是因恩许生的。}这些话是寓意说的。这两个妇女代表两个盟约：一个出自西乃山，生子为奴，那即是哈加尔；西乃山在阿拉伯，与现在的耶路撒冷相同，她和她的儿女同为奴隶。但那在上面的耶路撒冷是自由的，她是我们的母亲。}\par

9. 所以，这些恶人试图废除法律，却迫使我们赞同这些圣经。因为他们注意到经文所说：那些在法律之下的人为奴隶，他们不断重复那句话：\BibleQuote{你们这些靠法律成义的人，已与\Person{基督}隔绝，从恩宠上跌落了。}我们承认这一切都是真的，并且我们说，法律并非必要，除非对于那些奴役仍有益处的人；法律之所以有益地颁布，是因为那些无法被理性从罪恶中召回的人，需要被这样的法律——即被那些愚昧人能看见的惩罚的恐吓和恐怖——所约束。当\Person{基督}的恩宠释放我们时，它并不谴责那法律，而是最终邀请我们顺服于它的爱，而不是成为法律恐惧的奴隶。恩宠本身是白白恩赐，那些仍愿受法律束缚的人不明白恩宠是从天主而来。\Person{保禄}理所当然地斥责他们为不信者，因为他们不相信通过我们的主\Person{耶稣}，他们已从那被天主最公义的安排暂时置于其下的束缚中得到了释放。因此，同一位宗徒的话说：\BibleQuote{法律是我们的启蒙师，指向\Person{基督}。}所以，祂赐给人一个令人恐惧的启蒙师，后又赐给人一位令人爱的导师。然而，在法律的这些训令和命令中——这些如今基督徒不得使用的，如安息日、割损和祭献等——包含着如此伟大的奥秘，以致每个虔诚的人都能明白，没有什么比按文字（即按字面）理解其中任何内容更致命的了；也没有什么比在圣神中揭示它更有益的了。因此有这话：\BibleQuote{文字叫人死，圣神却叫人活。}因此又有：\BibleQuote{那同一帕子直到今还在朗读旧约时存留着，这帕子没有被除去，因为只在\Person{基督}内才被废除。}因为在\Person{基督}内被废除的，不是旧约，而是它的帕子；这样通过\Person{基督}，那在\Person{基督}之外晦暗隐藏的，就可以被理解，可以说是被揭示出来。因为同一位宗徒紧接着补充说：\BibleQuote{但你们若转向\Person{基督}，帕子就被除去。}因为他不是说法律被除去，或旧约被除去。所以，不是借着主的恩宠，好像那里隐藏着无用的东西而被除去；而是除去遮盖有用之物的遮盖物。凡真诚而虔诚地（而非无序和无耻地）寻求这些圣经意义的人，都这样对待他们；他们被仔细地展示事件顺序、言行原因，以及旧约与新约如此大的和谐，以至于没有一个点不一致；还有如此大的预象的秘密，以至于通过解释所引出的一切东西，都迫使他们承认自己是可悲的人，那些想在学会这些之前就谴责它们的人。\par

10. 但暂时撇开知识的深奥，以我应当对待密友的方式对待你，即按我自身的力量，而非我所惊叹的饱学之士的力量，我要谈谈人在阅读时犯错的三种情形。我将逐一说明。第一种情形：把虚假当作真实，而作者之意却非如此。第二种情形：虽不那么普遍，却同样有害，即把虚假当作真实，而所信与作者相同。第三种情形：从他人著作中领悟到某种真理，而作者本人却未理解。这种情形并非无益，若仔细思量，反而是阅读的全部果实。第一种情形的例子：比如有人因读了\Person{维吉尔}的诗句，便说并相信\Person{拉达曼图斯}在冥界审判死者的案件。此人误入歧途有二：一则相信了不该信之事，二则他所读的作者并不被认为相信此事。第二种情形可如此说明：若有人因\Person{卢克莱修}写道灵魂由原子构成，死后分解为原原子而消亡，便以为这是真实且应当相信的。此人同样不幸，若在如此重大的事上，将虚假当作确定之事，尽管\Person{卢克莱修}——其书使此人受骗——持此观点。此人虽确知作者之意，又有何益？他选择\Person{卢克莱修}并非为了逃避谬误，而是与他一同谬误。第三种情形的恰当例子：若有人读了\Person{伊壁鸠鲁}的某书中赞扬克制的段落，便断言他将\OriginalTerm{至善}{chief good}归于德行，因此不应受责。此人有何损于\Person{伊壁鸠鲁}的谬误？即便\Person{伊壁鸠鲁}相信肉体快乐是人的至善，此人并未屈从于如此卑劣有害的意见，他喜爱\Person{伊壁鸠鲁}，只因以为\Person{伊壁鸠鲁}未持不当之见。这种谬误不仅人之常情，也往往最合人性。试想，若有人告诉我一位我所爱的、已届成年的人，在众人面前说他如此喜爱童年与幼年，以至发誓希望过同样的生活，证据确凿，我若否认便属无耻；难道我若认为他此言意在表明他喜爱纯真与远离欲望纠缠的心境，并因此更爱他，这岂非应受责备？尽管他或许愚昧到喜爱孩童时期的游戏、食物与闲散？假设他在此消息传到我耳中后去世，我无法询问他解释其意；难道会有人无耻到因我通过所闻之词赞扬此人的心愿而恼怒吗？甚至一位公正的审判者或许也会称赞我的情感与心愿，因我喜爱纯真，且在可能作恶的情况下，作为人，对存疑之事宁愿善意揣测。\par

11. 既然如此，也请听同一圣经同样多的情形与区别。因为必然有同样多的情形出现在我们面前。或有人写了有益的东西，却未被某人有益地理解；或两者都无益；或读者有益地理解，而所读的作者却写了相反的东西。其中第一种我不责备，最后一种我不在意。因我不能责备一个无过而被误解的人；我也不能因某人未看到真理而烦恼，若我见读者并未受损。那么，有一种最受认可、似乎最纯净的情形，即所写的内容是好的，且被读者以善意接受。然而这仍可分为两类：因为它并未完全排除错误。通常发生的是：当作者持有善意的理解，读者也持有善意的理解，但彼此不同，且往往更好或更差，却仍有益。但当读者与作者持有相同理解，且此理解在生活方面全然正确时，真理便达到极致，也再无任何谬误可乘之隙。此种情形极为罕见，尤其是在所读内容极其晦涩时；依我之见，它不能被明确知晓，只能相信。因我凭何种证据能如此把握一位不在场或已故之人的意志，以至可以起誓？即使他当面受询，也可能有许多事——若他非恶人——他会极为谨慎地隐藏。但我认为，就了解事实而言，作者的品性无关紧要；然而，凡著作有益于人类及后世者，最应被公正地视为好人。\par

12. 因此，我希望他们告诉我，他们认为天主教会所谓的错误属于哪一种。如果是第一种，那是一个严重的指控；但不需要牵强的辩护，因为只要否认我们像攻击我们的人所设想的那样理解就够了。如果是第二种，指控同样严重；但他们将被同样的言论驳倒。如果是第三种，那就根本不是指控。接着，让我们考察\WorkTitle{圣经}本身。他们对所谓\WorkTitle{旧约}的书卷提出什么反对意见？是它们本身是好的，却被我们以错误的意思理解？但他们自己并不接受它们。或是它们既不好，也未被正确理解？但我们上面的辩护足以将他们从这一立场驱离。或者他们会说，虽然你们以好的意思理解，但它们本身是邪恶的？这岂不是等于开脱他们与之争论的活着的对手，而去指控早已死去、与他们并无争执的人？我确实相信那些人有益地将一切交付记忆，他们是伟大而神圣的。而那法律是按天主的命令和旨意颁布和制定的；对此，尽管我对这类书卷只有非常浅薄的认识，但我能轻易说服任何人，只要他以公正、不固执的心来听我；我会这样做，当你愿意以善意的耳朵和心灵倾听时；但这要等到我有能力的时候；而现在，无论此事情况如何，我未曾受骗，这难道还不够吗？\par

13. \Person{Honoratus}，我以我的良心和居于纯洁灵魂中的天主为证，我认为没有什么比天主教会以\WorkTitle{旧约}之名所保存的全部\WorkTitle{圣经}更明智、更纯洁、更虔诚的了。我知道你对此感到惊讶，因为我无法隐藏我们曾有不同的信念。然而，确实没有比这更鲁莽的事了（那时我们年轻时就是这样）：对于每一卷书，我们抛弃那些宣称持有并能够传授给弟子的解释者，而向那些人寻求其意义——我不知道出于什么原因，这些人对书的作者发起了最激烈的战争。有谁曾想过让\Person{亚里士多德}的敌人来解释其隐晦难懂的著作？且不提那些教学体系，其中读者或许可能犯错而不致亵渎。总之，有谁愿意在\Person{伊壁鸠鲁}的指导下阅读或学习\Person{阿基米德}的几何学著作？\Person{伊壁鸠鲁}曾以极大的固执反对这些著作，据我判断，他根本不懂它们。那些法律书卷最为清晰，而这些人却如同对公开之物徒劳无益地攻击，这算是什么？他们在我看来就像那个弱女子，正是这些男人常嘲笑的对象：她因一位摩尼教女信徒向她推崇太阳并劝她崇拜而愤怒，作为单纯而虔诚的人，她急忙跳起来，多次用脚踩踏太阳透过窗户照进光亮的那块地方，喊道：\Quote{看，我践踏太阳和你的天主！}完全是愚蠢和妇女式的行为，谁否认呢？但你觉得这些人不正是这样吗？他们对自己不理解的事物——无论是其缘由还是其性质——尽管这些事物对于理解者而言是精确而神圣的，他们却以大量的言词和辱骂进行攻击，以为自己在有所作为，因为无知者为他们喝彩。请相信我，这些\WorkTitle{圣经}中的一切，都是崇高而神圣的：其中全是真理，是最适合更新和滋养心灵的教学体系；而且安排得如此有分寸，以至于没有人不能从中汲取足够的养分，只要他怀着虔诚和孝心去汲取，正如真宗教所要求的。要向你们证明这一点，需要许多论证和更长的论述。因为我首先必须与你们讨论，使你们不憎恨作者本身；然后，使你们爱他们。而这一点我必须以其他方式论述，而非通过解释他们的含义和词句。原因如下：如果我们憎恨\Person{维吉尔}，更确切地说，如果我们不在理解他之前就爱他（凭我们先辈的称赞），那么关于他的无数问题（语法学家们常为此困扰不安）就永远无法令我们满足；我们也不会愿意倾听那些在赞扬他的同时解决这些问题的人；反而会偏向那些借此试图证明他犯错和糊涂的人。然而，如今许多人试图打开这些书卷，每个人按其能力以不同方式解释，我们更赞赏那些通过其解释使诗人显得更优秀的人——即使是那些不理解他的人也相信他不仅毫无过错，而且吟唱的无不配得赞美。因此，在某些细微的问题上，我们宁愿责怪失败的老师，他没有回答，也不会认为问题出在\Person{维吉尔}的沉默。而如果为了辩护自己，他想主张如此伟大作者的一个错误，他的学生几乎不会留下，即使他们已经付了学费。如果我们能对那些人表现出同样的善意，那将是多么伟大的事啊！自古以来那么长的时间已经证实，圣神藉他们说话。但是，我们这些最具理解力的年轻人、奇妙的理性探索者，至少没有展开这些书卷，没有寻求老师，没有责备自己的愚钝，最后，没有让自己的心向那些愿意使这类著作在全世界被如此长久阅读、保存和传布的人哪怕稍微屈服，却因那些对它们怀有敌意和仇恨者的言论，认为其中毫无可信之处——而那些人，在理性的虚假承诺下，迫使我们相信并珍视成千上万的寓言。\par

14. 但现在，我将继续我已开始的论述，若我能做到的话；我将这样与你们讨论：暂时不揭露天主教信仰，而是为了叫那些关心自己灵魂的人，能寻求其伟大奥迹，向他们显示神圣的果实和辨别真理的望德。凡寻求真宗教的人，无人怀疑他要么已相信有一个不朽的灵魂，使这宗教能裨益它，要么也愿意在这同一宗教中找到那事物本身。因此，一切宗教都是为了灵魂的缘故；因为，无论身体的本性如何，对于那灵魂已拥有使之受福的事物的人而言，身体，尤其是在死后，不会引起任何挂虑或焦虑。所以，为了灵魂的缘故，或唯此或首要地，真宗教——若真有这样的宗教——已被设立。但这灵魂（我将思考其原因，且承认此事极其隐晦）却仍是迷误和愚昧的，正如我们所见，直至它达到并察觉智慧，而或许这[智慧]就是真宗教。我岂是送你们去听神话？我岂是强迫你们轻率相信？我说，我们的灵魂纠缠并沉溺于谬误和愚昧中，寻求真理之路——若真有这样的路。若这并非你的情况，请宽恕我，我祈求，并与我分享你的智慧；但若你在我所说的身上认出你自己，那么我恳求，让我们一起寻求真理。\par

15. 假设我们尚未听说过任何宗教的导师。看，我们承担了一项新的事业和工作。我们必须寻找，我想，那些宣称此事业的人——若它确实存在。假设我们找到了持有不同意见的众人，他们因意见分歧而各自设法吸引人归向自己；但与此同时，有某些人因广为人知并几乎占据所有民族而显得卓越。他们是否持有真理，是一个重大问题；但难道我们不该先对他们进行充分的检验，以便当我们仍在迷误中（因为我们都是人）时，似乎是与人类本身一同迷误？\par

16. 但有人会说，真理存在于少数人那里；因此，若你知道它在谁那里，你就已知道它是什么了。我岂非刚才说过，我们是作为无知者在寻求它？但若你从真理本身的力量推断出拥有它的人很少，却不知他们是谁；那么，情况若如此：知道真理的人如此之少，以致他们以自己的权威掌握着群众，从而使这小部分人能得以解脱，并仿佛将自己挤入那些秘密中？难道我们看不见，达到最高雄辩术的人何其少，而全世界的修辞学校里却充满成群结队的青年？怎么，所有渴望成为优秀演说家的人，难道会因众多无知者的缘故而惊恐，认为应当把自己的努力放在\Person{凯基利乌斯}或\Person{埃鲁奇乌斯}的演说上，而非\Person{图利乌斯}的演说上？所有人都追求这些，它们因祖辈的权威而被确认。成群的无知者尝试学习同样的东西，而这同样被少数有学识者承认为值得学习之物；然而达到者极少，实践者更少，成为闻名者则少之又少。怎么，真宗教说不定也是如此？怎么，无知者成群结队前往教堂，但这并不证明因此没有人通过这些奥迹而成为成全的？然而，若学习雄辩术的人如同少数雄辩家那样稀少，我们的父母就绝不会相信我们应当被托付给这样的老师。那么，既然我们被一个多数——其中由无知者构成的部分为数众多——召唤去学习这些，以至于使我们爱上了少数人才能达到的事物；为何我们在宗教方面——它或许被我们以灵魂的巨大危险而轻视——却不愿处于同样的情形？因为，若最真实、最纯洁的天主敬拜，尽管只在少数人中被发现，却仍是与那些多数人——尽管他们沉溺于私欲，远离纯洁的理解——所一致的？（谁能否认这是可能发生的？）我问，若有人指责我们轻率和愚昧，说我们没有与那些教导它的人一起勤勉寻求那件我们极其渴望发现的事物，我们能回答什么呢？[我们是否要说：]我被人数吓退了？为什么在自由技艺的追求上——它们几乎不对此生带来什么益处；为什么在寻求钱财上？为什么在获得尊荣上？为什么，最后，在获得和保持健康上？还有，为什么在追求幸福生活上——尽管所有人都参与其中，但只有少数人出众——你却没有被任何人数吓退？\par

17. \Quote{但他们在那里似乎说了些荒谬的话。} 这是根据谁的断言？固然是敌人的断言，无论出于什么原因，无论什么理由，因为现在这不是问题所在，仍然是敌人。我亲自阅读时，发现确实如此。是这样吗？若没有接受过任何诗歌方面的教导，你不敢在没有老师的情况下尝试阅读\Person{Terentianus Maurus}；需要\Person{Asper}、\Person{Cornutus}、\Person{Donatus}和其他无数人，才能理解任何一位诗人，尽管他们的诗篇似乎甚至渴望剧院的掌声；你却对那些书籍——无论它们如何，但几乎全人类都普遍承认它们是神圣的、充满天主性的事物——在没有向导的情况下贸然阅读，并在没有老师的情况下敢于对它们发表意见；如果遇到任何似乎荒谬之处，你难道不该指责自己迟钝的心智、因这世界的败坏而衰败的心灵（所有愚昧之人都是如此），而非指责那些书籍吗？那些书籍也许根本不可能被这样的人理解！你应当立即寻找一位既虔诚又有学识的人，或者由许多人公认是如此的人，这样你就能从他的建议中获益，从他的学识中得到教诲。难道他不容易找到吗？应当努力寻找他。你居住的地方没有这样的人吗？有什么理由比旅行更能带来益处呢？难道他完全隐藏了，或者在大陆上不存在？应当渡海。如果渡海后在我们附近的任何地方都找不到他，你甚至应当前往那些书籍所述事件发生的地方。\newline{} 我们做了这类事情吗，\Person{Honoratus}？然而，也许是最神圣的宗教（因为我现在仍像在谈论一个可疑的问题），其意见如今已占据了整个世界，我们这些可怜的男孩却凭自己的判断和裁决定罪。如果那些圣经中冒犯某些无学识之人的事情，正是为了这个目的而放置在那里，即当我们读到那些违背常人（更不用说明智和圣洁的人）感受的事物时，我们会更加热切地寻求隐藏的含义。难道你未注意到，\WorkTitle{牧歌}中的\OriginalTerm{变童}{Catamite}——粗鲁的牧羊人为他泪流满面——人们试图解读，并断言那位少年\Person{Alexis}（据说\Person{柏拉图}也曾为他创作过一首爱歌）具有某种伟大的含义，但却是无学识之人无法理解的；而一位无论多么富有的诗人，似乎可以毫无亵渎地发表放荡的歌曲吗？\par

18. 然而，事实上，是否有任何法律的敕令、反对者的势力、献身者的卑劣品格、可耻的传闻、制度的创新或隐秘的职业，来阻止我们并禁止我们寻求？这些都没有。所有神律与人类法律都允许我们寻求天主教信仰；但持守并实践它，至少人类法律允许，即使我们处于错误中时对神律有所怀疑；没有敌人恐吓我们的软弱（尽管真理和灵魂的救恩，若在尽心寻求后未能在最安全之处找到，就应不惜任何风险去寻求）；所有等级和权力的阶层都最虔诚地服务于这神圣敬礼；宗教的名称最为尊贵、最为著名。我请问，有什么阻碍我们用虔诚而谨慎的探究去考察和讨论，这里是否有那样一种东西，它必须由少数人知道并完整地保存，尽管所有民族的善意和情感都赞同它？\par

19. 情况既然如此，假设正如我所说，我们现在是第一次寻求将我们的灵魂交付给哪一宗教，让它洁净并更新它们；无疑，我们必须从天主教会开始。因为现今基督徒的人数，已经超过犹太人和偶像崇拜者的总和。然而，在这些基督徒中，既然有几个异端，所有异端都希望显得像天主教徒，并把所有其他人都称为异端者，那么有一个天主教会，是所有人都承认的：如果你考虑整个世界，它在人数上更丰盈；但据知者证实，在真理上也比所有其他教会更纯洁。但真理的问题是另一回事；然而，对于寻求者来说，已经足够的是：有一个天主教会，不同的异端给它不同的名称，而他们自己却各自以自己特有的名称被称呼，他们不敢否认。由此，通过不受任何偏袒影响的仲裁者的判断，可以明白，那个所有异端都渴望的\Quote{天主教}名称，应当归属于谁。但为了避免有人认为这是一个过于啰嗦或不必要的讨论，至少这一点是确定的：人类法律本身在某种意义上也是基督徒的。我不希望从此事实形成任何偏见，但我认为这是探究最有利的开端。因为我们不必害怕，真正的天主崇拜，若本身没有力量，似乎需要那些它本当支持的人来支持；但无论如何，如果真理能在那里找到，那是最幸福的，因为在那里寻求和持守都最为安全；若不能，那么最终，无论冒什么风险，我们也必须去别处寻找。\par

20. 确立了这些原则后——\Quote{我认为这些原则如此公正，以至无论谁是我的对手，我理应在你们面前赢得这场诉讼}——我将尽己所能，向你们阐述我当年以那种精神寻求真宗教时所遵循的道路，那精神正如我现在所表明的，应以之寻求。因为离开你们、渡海之后，我迟疑不决，不知该持守什么、该放弃什么；这种迟疑自我们听说那人（你们知道，\Quote{如同从天上降下，要解释所有使我们困惑之事}）的到来之后，每日愈发增长，却发现他除了某种口才外，与常人无异；如今定居在\Place{意大利}，我反复思考、权衡，并非是否该继续留在那个我懊悔曾陷入的教派，而是如何找到真理——我对真理的渴慕，无人比你更清楚。我常觉得真理无法寻得，思绪的巨浪翻滚，倾向于赞同\OriginalTerm{学院派}{Academici}。又常竭尽所能，审视人类灵魂——如此有生命、如此有智慧、如此清明——以为真理并非隐藏，只是寻找真理的方式隐而不显，而这种方式必须来自某种神圣的权威。于是就需要探究那权威是什么，在如此众多的分歧中，每一派都承诺会传授它。于是我遇到了一座无法穿越的丛林，我非常不愿陷进去；在这些困扰中，我的心灵因渴望找到真理而无法安宁，不断翻腾。然而，我逐渐从那些我已决意离开的人中挣脱出来。但在如此大的危险中，除了以充满眼泪和悲伤的言词恳求天主眷顾帮助我之外，别无他法。我甘愿这样做；而如今\Place{米兰}主教的某些论述几乎促使我——并非没有希望——去探究许多关于\BibleBook{}{旧约}本身的事，你们知道，我们曾因它被错误地推荐给我们而视它为可诅咒的。我已决定在教会中成为一名慕道者（我的父母曾将我交付给这教会），直到我找到我所向往的，或说服自己无需再寻求为止。因此，如果有人能教导我，他必会发现我在一个非常关键的时刻，极为热忱且易于学习。如果你们看到自己长久以来也是如此，且怀着对灵魂同样的关怀，并且如今似乎觉得已经颠簸流离够了，希望结束这类劳苦，那么就跟随天主教教导的路径吧，这路径从\Person{基督}亲自经由宗徒们一直传到我们，并将在未来传到后世。\par

21. \Quote{这很可笑}，因为所有人都宣称持守并教导这一点；所有异端者都这样宣称，我不能否认；但他们是如此宣称，即向那些被他们引诱的人承诺，要就最晦涩的问题给出理由；并正因如此，他们主要指控天主教会，说来到她那里的人被要求相信；而他们却夸耀，自己不施加相信的轭，而是开启教导的源泉。你回答说：\Quote{还有什么比这更值得称赞的话呢？}并非如此。因为他们这样做，并非出于任何力量，而是为了以理性之名吸引群众；人的灵魂天性喜爱理性，却不顾自己的能力和健康状况，寻求健康者才应吃的食物（那食物只应托付给健康者），从而堕入欺骗者的毒药。因为真宗教，除非相信那些如果一个人行为端正、堪当相称之后才能达到并理解的事情，并且完全缺少某种有力的权威，否则绝不可能正确入门。\par

22. 但或许你要求就这一点给你一些理由，足以说服你不应当在信德之前先由理性教导。这很容易做到，只要你做一个公正的听者。然而，为了使它得体进行，我希望你仿佛回答我的问题；首先告诉我，你为什么认为人不应当相信。\Quote{因为在我看来，轻信（由此人们被称为轻信者）本身是一种过错；否则我们不会用它作为指责的用语。因为如果一个多疑的人犯了过错，在于他怀疑未确定的事；那么一个轻信的人更是如此，他与多疑者的区别在于，多疑者对未知之事留有一些怀疑，而轻信者则毫无怀疑。}我暂且接受这个看法和区别。但你知道，我们甚至通常用某种责备来称呼一个人为好奇者；但我们却用称赞来称呼他为勤学者。因此，请你注意，在你看来这两者之间有什么区别。\Quote{这当然是，虽然两者都被强烈的求知欲引导，但好奇者寻求与自己毫无关系的事，而勤学者则相反，寻求与自己相关的事。}但是，因为我们不否认一个人的妻子、儿女以及他们的健康与他相关；如果有人旅居在外，仔细询问所有来客他的妻子儿女如何、是否安好，他确实被强烈的求知欲引导，然而我们并不称这人为勤学者，尽管他非常想知道，而且这些事与他极其相关。因此你现在明白，勤学者的定义在这一方面出现缺陷：每个勤学者都想知道与自己相关的事，然而并非每个做这种事的人都可称为勤学者；而是指那些全心全意寻求与心智的自由修养和装饰相关之事的人。然而我们正确地称他为一个学习者，特别是如果我们加上他学习听什么。因为我们可以称他为自己家庭的勤学者，如果他只爱自己的家庭，但我们不加上某种限定，就不认为他配得上勤学者的通称。但一个渴望听到自己家人消息的人，我不会称他为勤听者，除非他因听到好消息而高兴，愿意一而再地听；但一个学习者，即使只听过一次。现在回到好奇者，告诉我，如果有人愿意听一些毫无益处的故事，即与他无关的事；而且并非以冒犯的方式经常听，而是很少且非常有节制，或在宴会上，或在某个聚会中，或任何集会中；他会显得好奇吗？\Quote{我认为不会；但无论如何，他肯定会显得关心他所愿意听的那件事。}因此，好奇者的定义也必须像勤学者的定义一样修正：那么请考虑先前的说法是否也需要修正。因为为什么一个有时怀疑某事的人不配被称为多疑者，而一个有时相信某事的人不配被称为轻信者？因此，正如一个学习任何事物的人与绝对勤学者之间有巨大差别；同样，关心某人与好奇者之间，相信某人与轻信者之间也有巨大差别。\par

23. 但你会说：\Quote{现在要考虑我们是否应当在宗教上相信。}因为，虽然我们承认相信是一回事，轻信是另一回事，但这并不意味着在宗教事务上相信就没有过错。因为如果相信和轻信两者都是过错，就像醉酒和酒徒两者都是过错一样，那又如何？现在，谁认为这是确定的，在我看来他就不能有朋友；因为，如果相信任何事情都是可耻的，那么要么他在相信朋友时行为可耻，要么他在任何事上都不相信朋友，我看不出他如何能称他为朋友或自称是朋友。在此你或许会说：\Quote{我承认在某些时候必须相信一些事；现在请说明，在宗教方面，为何在知道之前相信不是可耻的。}我会尽力说明。因此我问你，你认为更严重的过错是什么：将宗教传给不配的人，还是相信传道者所说的话？如果你不明白我称谁为不配的，我称那些以虚伪心怀接近的人。我想你同意，将任何神圣奥秘揭露给这样一个人，比相信宗教人士在宗教事宜上断言任何事情更应受责备。因为你不应作出其他回答。因此现在假设那位要把宗教传给你的人在面前，你如何能使他确信你是以真诚的心怀接近，并且就此事而言，在你里面没有欺诈或虚伪？你会说：\Quote{你自己的良心证明你毫无虚伪，用你所能的最强烈的言词来断言，但仍是言词。}因为你不能向人敞开你灵魂的隐藏之处，使人完全认识你。但如果他说：\Quote{看，我相信你，但你也相信我岂不是更公平？因为如果我持有任何真理，你将接受，我将给予，这是一个好处？}你会回答什么，除非是\Quote{你必须相信。}\par

24. 但是你说：\Quote{难道你不该给我一个理由，好使我能跟随那理由，无论它引向何处，而不带任何轻率吗？}或许应该如此；但既然这件事如此重大，你要藉理性达至对天主的认识，你认为所有人都能理解那些引导人的灵魂认识天主的理由吗？还是很多人，或是很少人？\Quote{很少人，}你说。\Quote{你相信自己在这些人之列吗？}你说，\Quote{这不该由我来回答。}因此，你认为他也该在这一点上相信你：而这他确实做了；只是你要记住，他已经两次相信了你说的不确定之事；而你却不愿相信他哪怕一次以虔诚精神劝告你。但假设如此，你以真诚的心态接近宗教，并且你是少数能够理解那些将神圣能力引入确定认识的理由的人之一；那么，你认为其他没有如此平和禀赋的人就该被拒绝宗教吗？或者你认为他们该被逐步引导，经由某些台阶到达那些最高的内在奥秘？你清楚哪一种更为虔诚。因为你不能认为，无论谁在渴望如此伟大之事时，都该以任何方式被遗弃或拒绝。但你难道不认为，除非他先相信他将达到他所立志的目标，并以恳求者的心态屈服，顺从某些伟大而必要的规诫，藉着某种生活方式彻底洁净自己的心灵，否则他绝不能达到真正纯净的事物吗？你当然这样认为。那么，那些（我已经相信你是其中之一）能够最轻易地藉确实理性接受神圣奥秘的人，如果他们像那些起初相信的人一样前来，这对他们会有任何妨碍吗？我认为不会。但是你说，\Quote{何必耽延他们？}因为他们虽然这样做绝不会伤害自己，却会因先例而伤害别人。因为几乎没有人对自己的能力有正确的认识：认识不足的人必须被激励；认识过高的人必须被抑制；这样不致使前者因失望而崩溃，也不致使后者因轻率而鲁莽。而这是容易做到的，如果连那些能飞翔的人（为免他们引诱任何人陷入危险）也被迫暂时走在别人也能安全行走的道路上。这就是真宗教的远见：这是天主的命令：这是从我们有福的祖先传承下来的，这是保存至今的：想要怀疑和推翻这一切，无异于寻找一条亵渎神圣的路通往真宗教。而这样做的人，即使他们所愿的被赐予了，也无法达到他们的目标。因为无论他们有怎样卓越的才智，若非天主临在，他们只能在地上爬行。但当那些追求天主的人顾念他们同为人类者时，祂就临在。没有比这一步更稳妥地朝向天堂了。就我而言，我无法抗拒这个推理，因为我怎能说我们不该相信任何没有确定认识的东西呢？然而，若没有相信某些无法以理性证明的事，就不可能有任何友谊；而且时常是身为奴隶的管家，无过错地被他们的主人信任。但在宗教中，还有什么比这更不公平：天主的仆人相信我们所承诺的真诚心志，而我们却不愿在他们命令我们任何事时相信他们？最后，还有什么方式比一个人藉着相信那些天主所指定、用以预先培育和陶冶心灵的事物，从而成为适合接受真理的人，更为有益健康？或者，如果你已经完全适合了，与其给自己带来危险并成为别人轻率的先例，不如稍微绕一下最安全的道路。\par

25. 因此，现在要考虑的是，我们应当以何种方式不跟随那些声称要凭理性引导我们的人。因为如何可以无罪地跟随那些命令我们相信的人，已经说过；但那些许诺理性的人，某些人认为他们不仅无罪，甚至可获称赞，但事实并非如此。因为有两种人（两类）在宗教上是可称赞的：一种是已经找到的人，我们应当判断他们是最有福的；另一种是极其热切并以正确方式寻求的人。因此，前者已经实在拥有，后者则行在路上，却是在那条必定抵达的路上。另有三种人全然应受斥责和厌恶。一种是那些持有意见的人，即那些自以为知道他们所不知道之事的人。另一种是那些确实知道自己不知道，却不如此寻求以致能够找到的人。第三种是那些既不自以为知道，也不愿寻求的人。还有三件事，在人的心灵中几乎彼此相邻，很值得区分：理解、相信、意见。如果这些各自考虑，第一种总是无罪的，第二种有时有罪，第三种从未无罪。因为对伟大、尊贵甚至神圣之事物的理解，是最有福的。但对不必要之事的理解并无伤害；也许学习本身是伤害，因为它占用了必要之事的时间。但就那些有害之事本身而言，不是理解，而是做或受它们，才是悲惨的。因为，假如有人理解如何无危险地杀死敌人，他并非因单纯的理解而有罪，而是因意愿；若意愿不存在，还有什么能称为更无辜呢？但相信只有在以下情形下才应受责备：要么对天主相信了不相称之事，要么对人的事过于轻信。在其他所有事上，若任何人相信某事，只要他明白自己不知道，就没有罪过。因为我相信从前极恶的谋反者被\Person{西塞罗}的德性处死；但这事我不仅不知道，而且我确知我绝不可能知道。但意见在两个方面极为可鄙：一是自认为已经知道的人无法学习（只要这学习是可能的）；二是鲁莽本身是不健全心灵的标志。因为即使有人假设他知道我关于\Person{西塞罗}所说的，（尽管这不妨碍他学习，因为这事本身无法被任何知识把握；）然而，由于他不明白有重大区别——某事是否被心灵的确定理性所把握（我们称之为理解），还是为了实际目的托付给传闻或记载，让后代去相信——他肯定犯错，而没有错误是不包含可鄙之事的。因此，我们所理解的，归于理性；我们所相信的，归于权威；我们所持意见的，归于错误。但每个理解的人也相信，每个持意见的人也相信；并非每个相信的人都理解，没有持意见的人理解。因此，如果这三件事关联到我们刚才提到的五种人，即可称赞的两种（我们首先列出）和有缺陷的三种，我们发现：第一种，即有福的人，相信真理本身；但第二种，即热切寻求并热爱真理的人，相信权威。在这两种人中，相信的行为是可称赞的。但在有缺陷的第一种人，即那些自以为知道他们所不知道之事的人中，有一种完全有缺陷的轻信。另外两种应受斥责的人什么也不相信，即那些寻求真理却绝望找不到的人，以及那些根本不寻求的人。而这只涉及属于任何教导体系的事。因为在生活的其他事务中，我完全不知道一个人如何能什么都不相信。尽管也有些人说在实际事务中他们遵循或然性，他们似乎更可能是不能知道而非不能相信。因为谁不相信他所赞同的呢？如果他们不赞同，他们所遵循的又怎能是或然的？因此，反对真理的人可能有两种：一种只攻击知识，不攻击信仰；另一种谴责两者。然而，我再次不知道在人类生活中是否能找到这样的人。这样说，是为了让我们理解，在持守信仰时，即使对尚未理解的事，我们也得以摆脱那些持有意见之人的鲁莽。因为那些说我们只相信我们所知道的人，是在提防那一个叫\Quote{持意见}的名称，这必须承认是可鄙且非常悲惨的；但若他们仔细考虑，自以为知道与出于某种权威相信他所不知道之事之间有极大区别，他们就一定会逃脱错误、不人道和骄傲的指责。\par

26. 我请问：如果未知之事不可信，那么子女如何能侍奉父母，并以彼此的爱去爱那些他们信以为父母的人呢？因为这事无论如何不能凭理性认识。但母亲的权威介入，使得我们相信父亲；而关于母亲，通常不是母亲本人被信，而是产婆、乳母、仆役。因为那母亲，她的儿子可能被偷换，另一个人被放在他的位置上，难道她就不会受骗而欺骗人吗？然而我们相信，并且毫无疑惑地相信，我们承认自己无法知道的事。因为若非如此，孝爱——人类最神圣的纽带——就会因极度的邪恶傲慢而被破坏，谁看不见这一点呢？因为即便是疯子，谁会认为那向自己信以为父母的人（即使他们并非真父母）尽本分的人该受责备呢？另一方面，谁不会判断那因害怕爱了冒充者而未能爱自己真父母的人应当被放逐呢？可以举出许多事例，表明如果我们决定不相信任何我们不能完全领会的东西，人类社会的任何部分都无法保持安全。\par

27. 但现在请听，我相信此时我能更容易地说服你。在宗教事务中，即对天主的敬拜和认识中，那些禁止我们相信，却急切宣称理性的人，更不应被跟随。因为无人怀疑所有人要么是愚人，要么是智者。但我现在所谓的智者，并非聪明而有天赋的人，而是那些在他们里面，有着尽可能多的人所能拥有的、对人和天主最确实的认识，以及适合那认识的生活和习俗的人；但其他所有人，无论其技能或无知，无论其生活行事是可赞或可责，我都将他们算在愚人之列。既然如此，稍有理智的人谁能不清醒地看出，愚人遵守智者的训诫，比凭自己的判断生活更有益、更健康呢？因为凡所行之事，若行得不对，就是罪；而非出自正确理性的事，无论如何也不能行得对。再者，正确理性即是德性本身。但除了智者的心灵，德性在谁那里可及呢？因此唯有智者不犯罪。因此，每个愚人犯罪，除非他在那些行动中服从了智者：因为所有这样的行动都出自正确理性，并且，可以说，愚人不应被视为自己行动的主人，他仿佛是智者的工具和侍从。所以，如果对所有人来说不犯罪比犯罪更好，那么无疑所有愚人若能成为智者的奴仆，他们会生活得更好。并且，若无人怀疑这在较小的事情上更好，例如买卖、耕种田地、娶妻、生养子女，最后在管理家产方面，那么在宗教方面更是如此。因为人类事务比神圣事务更容易分辨；而在所有更神圣、更崇高的事务中，我们对它们所应尽的服从和服务越大，犯罪就越邪恶、越危险。因此你从此看到，只要我们仍是愚人，且我们内心渴望一种卓越和虔诚的生活，除了寻求智者，别无他途；通过服从他们，我们既能减轻那统治我们的愚昧之感，并最终摆脱它。\par

28. 这里又出现一个非常困难的问题。因为我们愚人如何能找到一位智者呢？虽然几乎无人敢公开宣称，但大多数人都间接地声称这个名号：他们在智慧所包含的认识之事上如此彼此分歧，以至于要么他们中没有一个是智者，要么只有某一位是智者。但当愚人询问：\newline{}\Quote{那位智者是谁？}我完全看不出他如何能分辨和察觉。因为除非一个人已经认识那事物，而这些是其标记，否则他无法凭任何标记认出任何事物。但愚人不认识智慧。因为不像金银之类的东西，你可以看见它们却并不拥有它们；智慧不能被他人的心灵之眼所看见而不拥有它。因为凡通过身体感官接触的东西，都是外在地呈现给我们；因此我们能用眼睛感知属于他人的东西，即使我们自己并不拥有任何此类东西。但由理智所感知的东西是在心灵内部，拥有它无异于看见它。但愚人缺乏智慧，所以他不认识智慧。因为他不能用眼睛看见它：但他不能看见它却不拥有它，也不能拥有它却仍是愚人。因此他不认识它，并且只要他不认识它，他就不能在其他地方认出它。没有人，只要他是愚人，能以最确实的知识找出一位智者，通过服从他而摆脱如此大的愚昧之恶。\par

29. 因此，既然我们的探讨关乎宗教，如此巨大的困难唯有天主能予以补救；事实上，除非我们相信祂存在，且祂帮助人的心智，否则我们甚至不应该去寻求真正的宗教本身。因为我们如此努力追求的是什么呢？我们渴望达到什么呢？我们向往抵达何处？难道是我们不相信其存在或与我们有关的事物吗？没有什么比这种心态更荒谬的了。那么，当你不敢向我求取一件恩惠，或至少是无耻地敢于这样做时，你竟来要求发现宗教，而你认为天主既不存在，即使存在，也毫不关心我们？假如这件事如此重大，若不认真尽力寻求，就无法发现呢？假如发现本身极端困难，正是为了锻炼寻求者的心智，以接受那将要被发现的东西呢？因为还有什么比这光更令人愉悦和熟悉呢？然而人们经过长期黑暗后，竟不能承受和忍受它。还有什么比饮食更适合身体虚弱的人呢？然而我们看到，正在康复的人受到约束和限制，以免他们冒险像健康的人那样吃喝，从而因食物本身而重返厌弃食物的疾病。我指的是正在康复的人。至于那些病得很重的人，我们不也敦促他们吃点东西吗？他们若不相信自己能从这个疾病中痊愈，是绝不会如此痛苦地听从我们的。那么，你何时才会投身于充满痛苦和劳苦的探索呢？你何时才会有勇气为这件事承担如此大的关注和辛劳呢？当你并不相信你所寻求的东西存在时？因此，天主教教导体系的权威正确规定了：凡接近宗教的人，首先要被说服拥有信德。\par

30. 所以，那个异端者（因为我们讨论的是那些愿意被称为基督徒的人），我问你，他向我提出了什么理由？他凭什么理由叫我不要相信，仿佛相信是轻率之举？如果他叫我什么都不要相信，我就不相信在人间事务中存在任何真正的宗教；而我不相信存在的东西，我就不去寻找。但我想，他会在寻求时向我显示它：因为经上记载：\Quote{寻求的，就会找到。}因此，我不应该去他那里——他禁止我相信——除非我先相信某些东西。还有比这更疯狂的事吗？仅仅凭着那并非建立在知识上的信德（正是这信德引导我到他那里），我竟使他不悦？\par

31. 此外，所有异端者岂不都劝我们信从基督吗？他们难道还能更自相矛盾吗？在这件事上，可以双方面责问他们。首先，我们须问他们：他们素来应许的理性何在？对鲁莽的斥责何在？那自以为知道的知识何在？因为，若不加理由就相信是可耻的，那么你们在等待什么，又在忙碌什么，使我无理由地相信某一位，以便更容易被你们的理性引导？你们的理性又怎能将任何稳固的建筑建立在鲁莽的基础上？我是按那些因我们相信而不悦的人的口吻说话。因为我不仅认为在你们没有能力接受理性之前先相信，这是极有益于健康的，且借着信德本身彻底修养心灵以接纳真理的种子；而且，我认为这样一种事，若无它，患病的心灵便无法复原。既然这在你们看来是可嘲笑的、充满鲁莽的，那么你们还致力于使我们相信基督，岂不无耻？其次，我承认我已相信基督，并确信祂所说的都是真实的，尽管没有理性支持；异端者啊，这就是你首先要教导我的吗？容我稍作自省：既然我未曾亲眼见过基督自己——祂愿意向人显现时，据说曾被众人以肉眼看见——那么，我究竟是凭谁信了关于祂的事，以便我怀持此信德趋近你们？我发现我所信赖的，无非是众民与万邦的坚定意见和广泛流传的报道；而天主教会的奥迹，在任何时代、任何地方，都已拥有这些民众。那么，我为何不当优先向这些人审慎查明：基督究竟命令了什么？我早已因他们的权威而相信基督命令了有益之事。你难道能比我更好地解释祂的话吗？若由你推荐我这样相信，我原本连祂过去或现在的存在都不该相信。因此，如我所说，我是倚靠那由众多、一致与古老所巩固的报道而相信的。而你，既如此稀少，又如此喧嚣，如此新颖，谁都不会怀疑你提不出任何配得权威的东西。那么，这如此巨大的疯狂是什么？我相信他们，好让我相信基督，并从你们学习祂所说的？请问，为什么？即便他们失败，无法教导我任何事，我也更容易说服自己不该相信基督，而不是通过那些使我信祂的人以外的人去学习关于祂的任何事。哦，莫大的自信，毋宁说是荒谬！我教导你基督——你所信的那位——命令了什么。什么？如果我不信祂，你难道还能教导我关于祂的任何事吗？但他说，\Quote{你应当相信}。你的意思难道是我应当相信你，由你来引我信祂？不，他说，\Quote{因为我们是借着理性引导那些相信祂的人}。那么，我为何要信祂？因为报道已根深蒂固。是通过你们，还是通过别人？\Quote{通过别人}，他说。那么，我该相信他们，好让你来教导我？或许我该如此，若不是他们给我首要的嘱咐：\Quote{绝不可接近你们}；因为他们说\Quote{你们有致死的教义}。你会回答：\Quote{他们在撒谎}。那么，我该如何相信他们关于基督的事（他们未曾见过基督），却不相信他们关于你们的事（他们不愿看见你们）？\Quote{要相信圣经}，他说。但任何著作，若被当作新奇未闻之事提出，或由少数人推荐，而无理性证实，那么人们相信的不是那著作，而是提出它的人。因此，对于那些圣经，如果你是提出它们的人，既如此稀少又默默无闻，我不愿相信它们。同时，你也违背了你的承诺：强求信德而非给予理由。你又会引我回到众多与传闻。我请你，克制你的固执，以及这不知为何的、传播你名号的失控贪欲；倒要劝我去寻求这些群众中的首领，以全心勤苦向他们学习关于这些著作的事——若不是为了他们的存在，我根本不该知道我需要学习。但你该退回你的洞穴，不要以真理之名设下陷阱，试图从那些你自己也授予权威的人那里攫取真理。\par

32. 但若他们说，除非给我们确凿的理由，我们甚至不可相信基督，那他们便不是基督徒。因为这就是某些外邦人反对我们所说的话，确实愚昧，但与他们自己并不矛盾或相悖。但谁能容忍这些人自称属于基督，却争辩说除非他们向愚人提出关于天主最明了的理由，否则他们什么都不该相信？但我们看到，祂自己——就那部他们自己也相信的史书所教导的——先前所愿的，没有比祂应被相信更为强烈的事了；然而那些与祂打交道的人，还没有资格领受天主的奥秘。因为，那么多的伟大奇迹，祂自己也说，除了为使人相信祂，别无其他原因，还有什么其他目的呢？祂惯常用信德引导愚人，你们却用理智引导。祂惯常呼喊祂应被相信，你们却呼喊反对。祂惯常称赞那些相信祂的人，你们却责备他们。但除非祂将水变成酒——姑且不提其他（奇迹）——若人们跟随祂，祂却不做这些事，而只教导；那么要么我们必须不把那句话当回事：\Quote{你们要信天主，也要信我。}要么我们必须指责那人的轻率，他原不愿意让祂进入自己家，却相信只要祂一句话，他仆人的病就会痊愈。因此，祂给我们带来一剂良药，以医治我们彻底败坏的品行：借着奇迹为自己赢得权威，借着权威获得信德，借着信德招聚群众，借着群众拥有古老性，借着古老性巩固宗教；以致不仅行骗的异端者之极其愚蠢的新奇，而且外邦人暴力对抗的积年谬误，都无法从任何一面将它撕裂。\par

33. 因此，虽然我无力教导，但我仍不断劝告：鉴于许多人希望显得聪明，而要辨别他们是否是愚人也非易事，你当以全部热忱、全部祈祷，最后甚至以呻吟，或若可能，以眼泪，恳求天主救你脱离错误的邪恶；若你的心向往幸福生活。如果你以甘愿的心服从祂的命令——祂愿意这些命令由天主教会如此伟大的权威所确认——这事就会更容易成就。因为，既然智者与天主在心灵上如此结合，以致没有任何事物置于其间加以分离；因为天主是真理；除非他的心灵与真理接触，否则没有人是真正聪明的；我们不能否认，在人的愚昧与天主最纯全的真理之间，存在着人的智慧，作为某种中间之物。因为智者，就他被赋予的程度，效法天主；但对于愚人而言，没有比一个智慧的人更接近他、可供他有益地效法的了；并且，如前所述，用理智并不容易理解这智慧的人，所以必须让某些奇迹接近眼睛——愚人使用眼睛远比使用心灵更为便捷——这样，人们被权威所感动，他们的生活和习惯得以先被洁净，然后才能接受理智。因此，既然既需要效法人，又不可将我们的希望寄托于人，天主方面能做什么比这更充满慈爱和恩宠呢：那至纯、永恒、不变的天主智慧——我们理当依附于祂——竟然屈尊取了（人的）本性？这不仅是为了祂能做那些吸引我们跟随天主的事，也是为了祂能承受那些曾经阻止我们跟随天主的苦难。因为，除非一个人完全而完美地爱那最确实、至高的善，否则无人能达到它；而这决不会发生，只要身体的恶和命运的恶仍被畏惧；祂通过奇迹般的诞生和作为，使祂自己被爱；并通过死亡和复活，消除了恐惧。此外，在所有其他事情上——详细叙述会很长——祂显示自己如此，以致我们可以领悟到天主的仁慈能伸展到何等程度，以及人的软弱能被提升到何等程度。\par

34. 请相信我，这是最有益的权威，这是首先将我们的心从居住于地上抬举起来，这是从爱慕此世转向真天主。唯有权威能推动愚人急忙趋向智慧。只要我们尚不能理解纯正的真理，被权威误导固然可悲，但若不被感动则更为可悲。因为，如果天主眷顾不掌管人事，我们就无需为宗教操劳。然而，若一切事物的外在形式——我们深信它必然源自至真之美的泉源——以及某种内在的良心（我不知其为何物）仿佛在公共场所和私下里，劝勉所有较高层次的心思去寻求天主、事奉天主；那么，我们就不可完全放弃希望，认为同一天主自己已设立了某种权威，我们踏在其上，如同踏在一个稳固的阶梯上，便能被提升至天主面前。但这权威，撇开理智（我们常说，愚人很难理解纯正的理智），以两种方式感动我们：一部分通过奇迹，一部分通过跟随者的众多。这些对于智者而言，没有一样是必要的，谁能否认呢？但现在所要处理的事务是，使我们能够成为智者，即能够紧贴真理；而污秽的灵魂完全无法做到这一点：灵魂的污秽，简言之，就是爱慕任何非天主及灵魂本身的事物；人越是从这污秽中被洁净，就越容易看见真理。因此，为了洁净你的灵魂而渴望看见真理，殊不知灵魂正是为了看见而被洁净，这实在是颠倒和荒谬的。因此，对于无法看见真理的人，权威就在眼前，使他得以适合真理，并让自己被洁净；正如我上面所说，没有人怀疑这权威部分通过奇迹、部分通过群众而发挥作用。但我称那为奇迹：凡是出现困难或超乎惊讶者希望或能力之外的、不寻常之事。其中，对于民众，尤其是对于愚人，没有什么比贴近感官的事更合适了。但这些奇迹又分为两种：有些只引起惊讶，有些则同时带来极大的恩惠和善意。因为，若有人看见一个人飞翔，既然此事除了景观本身之外不给观看者带来任何益处，他就只是惊讶。但若有人身患重病且无希望，一听到命令便立刻痊愈，那么他对医治者的爱戴甚至超过对医治的惊讶。在真天主以真人显现于人的时代，这种事情发生得足够多了。病人被医治，麻风病人被洁净；跛子得以行走，瞎子得以看见，聋子得以听见。那时的人看见水变成酒，五千人用五个饼吃饱，海面被行走，死人复活：这些奇迹中，有些以更明显的恩惠服务于身体的益处，有些则以更隐蔽的标记服务于灵魂的益处，而所有奇迹都以对威能的见证服务于人的益处。就这样，天主的权威在那一时代，将世人漂泊的灵魂吸引到自己面前。你问：如今为何不再发生那些事？因为它们若不令人惊奇，就不会感动人；若成为平常，就不会令人惊奇。因为昼夜交替、天上事物的固定秩序、分为四季的年轮、树叶的飘落与复生、种子的无限力量、光的美、颜色、声音、滋味和气味的多样性——倘若有一个我们能与交谈的人初次看见并感知这些，他必会惊愕并陷入奇迹之中；然而我们却轻视这一切，并非因为它们易于理解（因为还有什么比这些的原因更晦涩呢？），而是因为它们不断接触我们的感官。因此，那些奇迹在非常合适的时机发生，为的是借着它们聚集并广传一群信徒，使权威有效地作用于习俗。\par

35. 然而，任何习惯都具有如此强大的力量，能占据人心，以至于即使是其中邪恶的——通常因过度贪欲而产生——我们也更容易指责和憎恨，而非放弃或改变。你以为这对人类的益处做得很少吗？不仅少数博学之士通过论证坚持，而且众多不同国家中无数未受教育的男女也都相信并宣告：我们不应崇拜任何地上的、火中的、最终与身体感官接触的事物为天主，而只应寻求藉着理智接近祂？他们禁食甚至到了只用面包和清水的最微薄食物，不仅终日禁食，而且连续数日；贞洁甚至达到了轻视婚姻和家庭；忍耐甚至到了轻视十字架和火焰；慷慨甚至到了将财产分给穷人；最后，对全世界的轻视甚至到了渴望死亡？少数人做这些事情，但更少人做得好且明智；然而整个民族赞同，民族聆听，民族赞成，最后，民族热爱。民族指责自己的软弱，因为他们无法做这些事情，而这并非没有将心灵提升至天主，也并非没有某些德行的火花。这是由天主眷顾所促成：藉着先知们的预言，藉着\Person{基督}的人性和教诲，藉着宗徒们的旅程，藉着殉道者们的侮辱、十字架、流血，藉着圣人们可赞美的生活，以及在这一切之中，按照时机的适宜，藉着与如此伟大事物和德行相称的奇迹。因此，当我们看到如此伟大的天主助佑、如此伟大的进展和成果，我们还会怀疑是否要藏身于那教会的怀抱中吗？她甚至从宗座开始，通过主教们的继承，得到了整个人类的承认（异端者徒然潜伏在她周围并被定罪，部分由于民众本身的判断，部分由于公会议的权威，部分也由于奇迹的庄严），持守着权威的顶峰。不愿将首位让给她，要么确实是不虔敬的极致，要么是轻率的傲慢。因为，如果没有确实的道路通向灵魂的智慧与健康，除非信仰为理性预备它们，那么，忘恩负义地拒绝神圣的帮助与援助，除了是希望抗拒以如此巨大辛劳装备起来的权威，还能是什么呢？并且，如果每一种教导体系，无论多么卑微浅显，为了被接受都需要一位导师或师傅，那么，在神圣奥秘的书卷方面，既不愿从解释它们的人那里学习，又想未经学习就谴责它们，还有什么比这更充满鲁莽的骄傲呢？\par

36. 因此，如果我们的推理或论述曾以任何方式触动了你，并且如果你，如我所信的，对自己有真正的关心，我愿你听从我，并以虔敬的信德、活泼的望德和纯朴的爱德，把自己托付给天主教信仰的良善导师；并不停地向天主本身祈祷，因为只有藉祂的良善我们被造，由祂的公义我们受罚，因祂的慈悲我们获得释放。这样，你既不会缺少最博学且真正基督徒的训诫和论著，也不会缺少书籍，甚至静思本身，藉此你可以轻易寻获你所寻求的。因为你要彻底抛弃那些饶舌而可怜的人（因为我能使用什么更温和的名称呢？），当他们过度寻求恶的来源时，所发现的除了恶以外没有别的。关于这个问题，他们常常激起听者去探究；但一旦被激起后，他们所教导的内容竟是：好永远沉睡，也不愿如此清醒。因为他们以疯狂代替昏睡，在这两种通常致命的疾病之间，仍有这样的区别：昏睡者死亡而不对他人施暴；但疯狂者却使许多健康的人，特别是那些想要帮助他的人，有理由恐惧。因为天主既不是恶的作者，也从未后悔祂所做的任何事，也不被任何灵魂情感的暴风雨所困扰，祂的王国也不是地上的一小部分：祂从不认可或命令任何罪恶或不义，祂从不说谎。因为这些以及诸如此类的说法，从前曾激动我们，当他们以此发动巨大而威胁性的攻击，并指控这正是旧约的教导体系时——这是极其虚假的。因此我承认，他们在谴责这些事上是正确的。那么我学到了什么？你以为如何？无非是，当这些被谴责时，天主教的教导体系并未受到谴责。因此，我在他们中所学到的真实之事，我持守；我以为真实而实为虚假的，我拒绝。但天主教会也教导了我许多其他事物，那些有躯体无血肉、但思想粗劣的人不能企及；也就是说，天主不是有形体的，祂的任何部分都不能被肉眼感知，祂的\OriginalTerm{实体或本性}{Substance or Nature}绝不可能遭受暴力或变化，也非组合或形成；如果你允许我这些（因为我们对天主不能有别的想法），那么他们的一切计谋都被推翻了。但天主既未生也未造恶，也不存在任何天主未生或未造的本性或实体，然而祂却使我们脱离邪恶，这是如何做到的，已经由如此必然的理由所证明，以至于完全不容置疑；尤其是对你以及像你这样的人；也就是说，如果在良好的性情上再加上虔诚和一定的心平气和，没有这些，对于如此重大的事就完全无法理解。而这里没有关于烟雾的传闻，也没有我不知道的某个\OriginalTerm{波斯的虚妄传说}{Persian vain fable}，只需侧耳倾听即可，而这样的灵魂并不敏锐，而是完全幼稚。事实远非如此，远非\OriginalTerm{摩尼教徒}{Manichees}痴人说梦那样。但既然我们这番论述已远超我的预料，那就让我们在此结束这卷书；在此书中我希望你记住，我尚未开始驳斥摩尼教徒，也尚未攻击那些胡言乱语；而且我也没有阐述任何关于天主教会本身的伟大之事，而只是希望，如果可能，从你心中铲除一个关于真正基督徒的错误观念，这观念是恶意或无知地灌输给我们的，并激发你去学习某些伟大而神圣的事。所以，让这卷书就这样吧；但当你的灵魂更加平静时，对于余下的事，我或许会更加乐意。\par

\SourceRecord{Translated by C.L. Cornish.；Nicene and Post-Nicene Fathers, First Series；Vol. 3.；Edited by Philip Schaff.；Buffalo, NY: Christian Literature Publishing Co.,；1887.；<http://www.newadvent.org/fathers/1306.htm>.}
